\documentclass[11pt]{article}

\usepackage[a4paper,margin=1in]{geometry}
\usepackage{amsmath,amssymb,amsthm}
\usepackage{enumitem}
\usepackage{xcolor}
\usepackage{fancyhdr}
\usepackage{needspace}
\usepackage{physics}

\definecolor{courseblue}{RGB}{0,0,255}

\setlength{\parindent}{0pt}
\setlength{\parskip}{0.35em}
\setlength{\headheight}{14pt}

\pagestyle{fancy}
\fancyhf{}
\fancyfoot[C]{\thepage}
\renewcommand{\headrulewidth}{0pt}

\renewcommand{\ip}[2]{\left\langle #1, #2 \right\rangle}

\usepackage{bm}
\renewcommand{\vec}[1]{\bm{#1}} % Converts all \vec{} commands to bold

\DeclareMathOperator{\dist}{dist}
\DeclareMathOperator{\open}{open}
\DeclareMathOperator{\closed}{closed}
\newtheorem*{lemma}{Lemma}
\newcounter{problem}
\newenvironment{problem}{%
    \newpage
    \refstepcounter{problem}%
    \Needspace{7\baselineskip}%
    \par\vspace{0.7em}%
    \noindent\textcolor{courseblue}{\rule{\textwidth}{0.4pt}}%
    \par\vspace{0.5em}%
    \begin{list}{}{%
        \setlength{\leftmargin}{2.3em}%
        \setlength{\labelwidth}{1.8em}%
        \setlength{\labelsep}{0.5em}%
        \setlength{\topsep}{0pt}%
        \setlength{\partopsep}{0pt}%
    }%
    \item[{[\theproblem]}]
}{%
    \end{list}
    \par\vspace{0.5em}
}

\newenvironment{parts}{%
    \begin{enumerate}[label={[\alph*]},leftmargin=2.2em,itemsep=0.65em,topsep=0.65em]
}{%
    \end{enumerate}
}

% Type each solution between \begin{answer} and \end{answer}.
\newenvironment{answer}{%
    \Needspace{4\baselineskip}%
    \par\vspace{0.35em}
    \noindent\textbf{Answer:}\par
}{%
    \par\vspace{0.25em}
}

\begin{document}

\begin{center}
    {\color{courseblue}\Large\textbf{AMATH/PMATH 331: Applied Real Analysis}}
    \hfill
    {\color{courseblue}\Large\textbf{Fall 2026}}

    \vspace{1.2em}
    {\color{courseblue}\Large
    Assignment 4 Submission; due Tuesday, November 10, 2026}

    \vspace{1.2em}
    
    \textbf{Name:} Vlad Paun

    \textbf{Student number:} 21291453
\end{center}

\begin{problem}
Let $A\subseteq\mathbb{R}^n$ be a non-empty subset, and let
$\operatorname{Map}(A,\mathbb{R}^m)$ be the space of maps from $A$ to $\mathbb{R}^m$.
\begin{parts}
    \item Let $(f_k)$ be a sequence in $\operatorname{Map}(A,\mathbb{R}^m)$
    which converges uniformly to a \emph{bounded} function $f$, so
    $f\in B(A,\mathbb{R}^m)$. Prove that the $f_k$ are \emph{eventually bounded}.
    That is, prove that there exists $N\in\mathbb{N}$, such that $k\geq N$
    implies $f_k\in B(A,\mathbb{R}^m)$.
    \begin{answer}
        \begin{proof}
            Since $f_k$ converges uniformly to $f$:
            \[\forall \varepsilon > 0:\ \exists N\in\mathbb{N}:\ \forall k \geq N:\ \norm{f_k(x) - f(x)} < \varepsilon \quad \forall x\in A\]
            So take $\varepsilon = 1$ and the respective $N$.
            Also, since $f$ is bounded, take $M$ such that $\forall x\in A:\ \norm{f(x)} \leq M$

            Now,
            \begin{align*}
                \norm{f_k(x)} & = \norm{f_k(x) - f(x) + f(x)} & \forall x\in A\\
                & \leq \norm{f_k - f(x)} + \norm{f(x)} & \forall x\in A\\
                & < \varepsilon + M = 1 + M & \forall k \geq N, \forall x\in A
            \end{align*}

            Therefore, $\forall k \geq N:\ f_k\in B(A,\mathbb{R}^m)$
        \end{proof}
    \end{answer}

    % The source uses R^n here and R^m for the limits; preserved as printed.
    \item Suppose $(f_k)$ and $(g_k)$ are two sequences in
    $\operatorname{Map}(A,\mathbb{R}^n)$ which converge uniformly to
    \emph{bounded} functions $f,g\in B(A,\mathbb{R}^m)$. Prove that the
    sequence $(\langle f_k,g_k\rangle)_{k\in\mathbb{N}}$ in
    $\operatorname{Map}(A,\mathbb{R})$ converges uniformly to
    $f\cdot g\in\operatorname{Map}(A,\mathbb{R})$, where
    $\langle\cdot,\cdot\rangle$ is the dot product in $\mathbb{R}^m$.
    \begin{answer}
        \begin{proof}
            We have to show
            \[\forall \varepsilon > 0:\ \exists N\in\mathbb{N}:\ \forall k \geq N:\ \abs{\ip{f_k}{g_k}(x) - \ip{f}{g}(x)} < \varepsilon \quad \forall x\in A\]

            By $1[a]$, $g_k$ is eventually bounded, so let $N_0$ such that $\forall k \geq N_0:\ \norm{g_k(x)} \leq M_0$.
            
            Let $M_f,M_g$ be the bounds of $f,g$.

            Let arbitrary $\varepsilon > 0$. Let $\varepsilon' = \frac{\varepsilon}{M_0 + M_f}$. Let $N_f,N_g$ such that:
            \[\forall k\geq N_f:\ \norm{f_k(x) - f(x)} < \varepsilon' \quad \forall x\in A\]
            \[\forall k\geq N_g:\ \norm{g_k(x) - g(x)} < \varepsilon' \quad \forall x\in A\]

            Take $N = \max(N_f,N_g,N_0)$ and take $k\geq N$.

            Now,
            \begin{align*}
                \abs{\ip{f_k}{g_k}(x) - \ip{f}{g}(x)} & = \abs{\ip{f_k}{g_k}(x) - \ip{f}{g_k}(x) + \ip{f}{g_k}(x) - \ip{f}{g}(x)}\\
                & = \abs{\ip{f_k - f}{g_k}(x) + \ip{f}{g_k - g}(x)}\\
                & \leq \abs{\ip{f_k - f}{g_k}(x)} + \abs{\ip{f}{g_k - g}(x)}\\
                & \leq \norm{(f_k(x) - f(x))}\norm{g_k(x)} + \norm{f(x)}\norm{(g_k(x) - g(x))}\\
                & < \varepsilon'\norm{g_k(x)} + \varepsilon'\norm{f(x)}\\
                & < \varepsilon'(M_0 + M_f) = \frac{\varepsilon}{M_0 + M_f}(M_0 + M_f) = \varepsilon
            \end{align*}
        \end{proof}
    \end{answer}

    \item Show that the assumption of boundedness of the limits in part [b]
    is essential as follows. Show that the sequence $(f_k)$ in
    $M(\mathbb{R},\mathbb{R})$ given by $f_k(x)=x+\frac{1}{k}$ converges
    uniformly on $\mathbb{R}$ to the unbounded function $f(x)=x$, but that
    $f_k^2\in M(\mathbb{R},\mathbb{R})$ does not converge uniformly to $f^2$.
    \begin{answer}
        \begin{proof}
            Let arbitrary $\varepsilon > 0$. Take $N > \frac{2}{\varepsilon}$
            \begin{align*}
                \abs{f_k(x) - f(x)} & = \abs{x + \frac1k - x} = \abs{\frac1k} \\
                & \leq \frac1N < \frac\varepsilon2 < \varepsilon & \forall k\geq N,\forall x\in \mathbb{R}
            \end{align*}
            So $f_k$ converges uniformly to $f$.

            Now, We have to show
            \[\exists \varepsilon > 0:\ \forall N\in\mathbb{N}:\ \exists k \geq N:\ \exists x\in\mathbb{R}:\ \abs{f_k^2(x) - f^2(x)} \geq \varepsilon \]
            
            Take $\varepsilon = 1$. Let arbitrary $N\in\mathbb{N}$. Take $k = N$. Take $x = k$
            
            \begin{align*}
                \abs{f_k^2(x) - f^2(x)} & = \abs{x^2 + \frac{2x}{k} + \frac{1}{k^2} - x^2} = \abs{\frac{2x}{k} + \frac{1}{k^2}}\\
                & = \abs{2 + \frac{1}{N^2}} > 2 > 1 = \varepsilon
            \end{align*}
            So $f_k^2$ does not converge uniformly to $f^2$.
        \end{proof}
    \end{answer}
\end{parts}
\end{problem}

\begin{problem}
Let $f\in\operatorname{Map}(\mathbb{R}^n,\mathbb{R}^m)$ be a map, and
suppose that $f$ is \emph{continuous} at $\mathbf{x}_0\in\mathbb{R}^n$.

Define a sequence $(f_k)$ in $\operatorname{Map}(\mathbb{R}^n,\mathbb{R}^m)$
by $f_k(\mathbf{x})=f(\mathbf{x}_0+\frac{1}{k}\mathbf{x})$.
Let $g\in\operatorname{Map}(\mathbb{R}^n,\mathbb{R}^m)$ be the
\emph{constant map} $g(\mathbf{x})=f(\mathbf{x}_0)$ for all
$\mathbf{x}\in\mathbb{R}^n$.

Prove that, for any $R>0$, the sequence $(f_k)$ converges uniformly to $g$
on the set $B_R(\mathbf{0})$.
\begin{answer}
\end{answer}
\end{problem}

\begin{problem}
For $k\in\mathbb{N}$, define $f_k:[-1,1]\to\mathbb{R}$ by
$f_k(x)=(1-x^2)^k$.
\begin{parts}
    \item For any $\delta\in(0,1)$, prove that the sequence of functions
    $(f_k)$ \emph{converge uniformly} on the subset
    $[-1,-\delta]\cup[\delta,1]$ of $\mathbb{R}$.
    \begin{answer}
    \end{answer}

    \item Prove that the sequence of functions $(f_k)$ \emph{does not
    converge uniformly} on the subset $[-1,1]$ of $\mathbb{R}$.
    \begin{answer}
    \end{answer}
\end{parts}
\end{problem}

\begin{problem}
Let $A\subseteq\mathbb{R}^n$ be non-empty. Let $(f_k)$ be a sequence in
$\operatorname{Map}(A,\mathbb{R}^m)$ which converges \emph{uniformly} to a
limit $f\in\operatorname{Map}(A,\mathbb{R}^m)$. We proved in class that if
every $f_k$ is continuous on $A$, then $f$ is also continuous on $A$.
\begin{parts}
    \item Adapt that proof to show that if every $f_k$ is \emph{uniformly}
    continuous on $A$, then $f$ is also \emph{uniformly} continuous on $A$.
    \begin{answer}
    \end{answer}

    \item Explain why we do not need to do any extra work if $A$ is compact.
    That is, why does part [a] follow immediately if $A$ is compact?
    \begin{answer}
    \end{answer}
\end{parts}
\end{problem}

\begin{problem}
COMING SOON
\begin{answer}
\end{answer}
\end{problem}

\begin{problem}
COMING SOON
\begin{answer}
\end{answer}
\end{problem}

\begin{problem}
COMING SOON
\begin{answer}
\end{answer}
\end{problem}

\begin{problem}
COMING SOON
\begin{answer}
\end{answer}
\end{problem}

\begin{problem}
COMING SOON
\begin{answer}
\end{answer}
\end{problem}

\begin{problem}
COMING SOON
\begin{answer}
\end{answer}
\end{problem}

\end{document}
